% PTE Core 模考 1B —— 听力 错题与详解·整理卷  版本 001（自包含）
% 编译: xelatex 本文件.tex   （需 Noto CJK + TeX Gyre Heros 字体）
% 本版本无外部图片，单文件即可复现。
\documentclass[11pt]{article}

% --------- 样式（原 ptestyle.sty，已内联）---------
% ============================================================
%  ptestyle.sty  —  PTE 模考错题整理卷 共享样式
%  编译引擎: XeLaTeX  (中英混排, 需要 Noto CJK 字体)
% ============================================================

\usepackage[a4paper,margin=2cm,headsep=4mm]{geometry}
\usepackage{fontspec}
\usepackage{xeCJK}
\usepackage[table]{xcolor}
\usepackage[normalem]{ulem}        % \sout 删除线
\usepackage{tabularx}
\usepackage{array}
\usepackage{ragged2e}
\usepackage{enumitem}
\usepackage[most]{tcolorbox}
\usepackage{fancyhdr}
\usepackage{graphicx}
\usepackage{pifont}                % \ding{51}=✓ \ding{55}=✗（阅读对错标注）
\usepackage{needspace}
\usepackage{tikz}

% ---------- 字体 ----------
\setmainfont{TeX Gyre Heros}            % 拉丁: Helvetica 风格(纤细清晰)
\setCJKmainfont{Noto Sans CJK SC}
\setCJKsansfont{Noto Sans CJK SC}
\newfontfamily\cjkserif{Noto Serif CJK SC}
\linespread{1.18}
\setlength{\parindent}{0pt}
\setlength{\parskip}{4pt}

% ---------- 配色 ----------
\definecolor{cWrong}{HTML}{D32F2F}     % 红  = 修改 / 错误
\definecolor{cFix}{HTML}{1B7F3B}       % 绿  = 正确建议
\definecolor{cDelete}{HTML}{1565C0}    % 蓝  = 删除
\definecolor{cInsert}{HTML}{7B1FA2}    % 紫  = 插入
\definecolor{cPartial}{HTML}{E07B00}   % 橙  = 部分得分
\definecolor{cNote}{HTML}{0277BD}      % 蓝  = 注解
\definecolor{cAccent}{HTML}{16A085}    % 主题青绿(平台主色)
\definecolor{cWrongBg}{HTML}{FBE0E0}
\definecolor{cDelBg}{HTML}{DCE8FB}
\definecolor{cInsBg}{HTML}{EEE0F5}
\definecolor{cHead}{HTML}{20303A}
\definecolor{cGrayBox}{HTML}{F4F6F7}
\definecolor{cModelBg}{HTML}{EAF7EF}
\definecolor{cYouBg}{HTML}{FFF7F5}
\definecolor{cNoteBg}{HTML}{EAF3FA}
\definecolor{cRule}{HTML}{D7DEE2}

% ---------- 行内批改宏 (复刻平台标注) ----------
% \fix{原词}{正确}  红色删除线 + 绿色上标改正  (修改/拼写/空格)
\newcommand{\fix}[2]{%
  \colorbox{cWrongBg}{\textcolor{cWrong}{\sout{#1}}}%
  \,\textsuperscript{\textcolor{cFix}{\footnotesize #2}}}
% \del{原词}  蓝色删除线  (应删除)
\newcommand{\del}[1]{%
  \colorbox{cDelBg}{\textcolor{cDelete}{\sout{#1}}}%
  \,\textsuperscript{\textcolor{cDelete}{\footnotesize 删}}}
% \ins{词}  紫色下划线  (应插入)
\newcommand{\ins}[1]{%
  \textsuperscript{\textcolor{cInsert}{\footnotesize $\wedge$}}%
  \colorbox{cInsBg}{\textcolor{cInsert}{#1}}}
% \miss{词}  灰色删除线  (口语漏读 / 未作答)
\newcommand{\miss}[1]{\textcolor{gray}{\sout{#1}}}
% \add{原词}{改正}  橙色波浪线 + "补"标签  (平台漏标、本卷补标的错误)
\definecolor{cAdd}{HTML}{E65100}
\newcommand{\add}[2]{%
  \textcolor{cAdd}{\uwave{#1}}%
  \,\textsuperscript{\textcolor{cAdd}{\footnotesize 补：#2}}}
% 高亮(口语发音): 好/中/差
\newcommand{\spk}[2]{\textcolor{#1}{#2}}

% ---------- 评分单元格着色 ----------
\newcommand{\sfull}[1]{\textcolor{cFix}{\bfseries #1}}     % 满分
\newcommand{\spart}[1]{\textcolor{cPartial}{\bfseries #1}} % 部分
\newcommand{\szero}[1]{\textcolor{cWrong}{\bfseries #1}}   % 零/低分

% ---------- 分数小标签 ----------
\newtcbox{\chip}[1][cAccent]{on line, boxsep=2pt, left=4pt, right=4pt, top=1pt,
  bottom=1pt, arc=3pt, colback=#1, colframe=#1, coltext=white, nobeforeafter,
  fontupper=\bfseries\footnotesize}

% ---------- 题块小标题 ----------
\newcommand{\ptesub}[1]{%
  \par\needspace{2\baselineskip}\smallskip
  {\color{cAccent}\rule[-0.2ex]{3pt}{1.05em}}\hspace{6pt}%
  {\bfseries\color{cHead}#1}\par\nopagebreak\smallskip}

% ---------- 题块外框 ----------
% \begin{qblock}{标号·题型}{右上信息(得分/用时)} ... \end{qblock}
\newtcolorbox{qblock}[2]{
  breakable, enhanced, arc=2.5pt,
  colback=white, colframe=cRule, boxrule=0.6pt,
  left=11pt, right=11pt, top=8pt, bottom=10pt,
  toptitle=3pt, bottomtitle=3pt, lefttitle=11pt, righttitle=11pt,
  colbacktitle=cHead, coltitle=white, fonttitle=\bfseries,
  title={#1\hfill\normalfont\bfseries\small #2}}

% ---------- 文本块: 题干 / 范文 / 你的作答 / 建议 ----------
\newtcolorbox{promptbox}{enhanced, breakable, colback=cGrayBox, colframe=cGrayBox,
  arc=2pt, left=8pt, right=8pt, top=6pt, bottom=6pt, boxrule=0pt}
\newtcolorbox{modelbox}{enhanced, breakable, colback=cModelBg, colframe=cModelBg,
  borderline west={3pt}{0pt}{cFix}, arc=1pt, left=8pt, right=8pt, top=6pt, bottom=6pt, boxrule=0pt}
\newtcolorbox{youbox}{enhanced, breakable, colback=cYouBg, colframe=cYouBg,
  borderline west={3pt}{0pt}{cWrong}, arc=1pt, left=8pt, right=8pt, top=6pt, bottom=6pt, boxrule=0pt}
\newtcolorbox{notebox}{enhanced, breakable, colback=cNoteBg, colframe=cNoteBg,
  borderline west={3pt}{0pt}{cNote}, arc=1pt, left=8pt, right=8pt, top=6pt, bottom=6pt, boxrule=0pt}

% ---------- 评分表 ----------
% 用法见正文; 这里给表头宏
\newcommand{\scorehead}{%
  \rowcolors{2}{cGrayBox}{white}
  \arrayrulecolor{cRule}}

% ---------- 章节大标题 ----------
\newcommand{\ptesection}[3]{% #1=中文名 #2=英文名 #3=该项分数
  \needspace{6\baselineskip}
  \begin{tcolorbox}[enhanced, colback=cAccent, colframe=cAccent, sharp corners,
     left=14pt, right=14pt, top=10pt, bottom=10pt, boxrule=0pt]
    {\fontsize{20}{24}\selectfont\bfseries\color{white} #1\ \ }%
    {\large\color{white!85} #2}\hfill
    {\large\color{white}本项得分\ \ }{\fontsize{22}{24}\selectfont\bfseries\color{white} #3}%
    \ {\color{white!85}/ 90}
  \end{tcolorbox}}

% ---------- 题型小节分隔 (口语 RA/RS/DI/RTS/ASQ 等) ----------
\newcommand{\ptetask}[2]{%
  \needspace{4\baselineskip}\medskip
  \begin{tcolorbox}[enhanced, colback=cHead, colframe=cHead, sharp corners,
     left=12pt, right=12pt, top=5pt, bottom=5pt, boxrule=0pt]
    {\large\bfseries\color{white} #1}\ \ {\color{white!80} #2}
  \end{tcolorbox}\smallskip}

% ---------- 口语 AI 语音识别着色 (复刻平台: 优/良/差 + 停顿) ----------
\definecolor{cSpkGood}{HTML}{2E9E5B}   % 优 绿
\definecolor{cSpkOk}{HTML}{C77F00}     % 良 黄/琥珀
\definecolor{cSpkBad}{HTML}{E0352B}    % 差 红
\newcommand{\sg}[1]{\textcolor{cSpkGood}{#1}}   % 优
\newcommand{\sy}[1]{\textcolor{cSpkOk}{#1}}     % 良
\newcommand{\sr}[1]{\textcolor{cSpkBad}{#1}}    % 差
\newcommand{\smiss}[1]{\textcolor{cSpkBad}{\sout{#1}}}  % 漏读 (红删除线)
\newcommand{\pse}{{\color{gray}\,/\,}}          % 停顿
\newcommand{\asrlegend}{{\footnotesize\color{gray}AI 语音识别\quad
  {\color{cSpkGood}\textbullet}\,优\ \ {\color{cSpkOk}\textbullet}\,良\ \
  {\color{cSpkBad}\textbullet}\,差\ \ {\color{gray}/}\,停顿\ \
  {\color{cSpkBad}\sout{词}}\,漏读}}
\newtcolorbox{asrbox}{enhanced, breakable, colback=white, colframe=cRule,
  boxrule=0.6pt, arc=1pt, left=8pt, right=8pt, top=6pt, bottom=6pt}

% ---------- 中文翻译行 ----------
\newcommand{\zh}[1]{{\footnotesize\textcolor{cNote}{\textbf{译}}\ \textcolor{gray}{#1}}}

% ---------- 阅读填空 / 选择 对错标注 ----------
\newcommand{\rok}[1]{{\color{cFix}\ding{51}\,#1}}                 % 选对(绿勾 + 词)
\newcommand{\rno}[2]{{\color{cWrong}\ding{55}\,\sout{#1}}\,{\footnotesize\color{cFix}(答案:\,#2)}}  % 选错: 你的(红删) + (答案:正确)
\newcommand{\rblank}[1]{{\color{cFix}\underline{\,#1\,}}}         % 直接给空的正确答案(无你的作答时)

% ---------- 听力专用 ----------
\newcommand{\hiw}[2]{{\color{cWrong}\uline{#1}}\,{\footnotesize\color{cFix}(听:#2)}} % HIW: 屏幕错词(红下划线)+ 实际读音
\newcommand{\lhit}[1]{{\color{cFix}#1}}                           % WFD: 命中的词(绿)
\newcommand{\lmiss}[1]{{\color{cWrong}\sout{#1}}}                 % WFD: 多余/错误的词(红删)

% ---------- 颜色图例 ----------
\newcommand{\ptelegend}{%
  \begin{tcolorbox}[enhanced, colback=cGrayBox, colframe=cRule, boxrule=0.5pt,
    arc=2pt, left=10pt, right=10pt, top=6pt, bottom=6pt]
  {\bfseries\color{cHead}颜色说明\quad}
  \fix{改}{正}~修改/拼写\quad
  \del{删}~应删除\quad
  \ins{插}~应插入\quad
  \add{补}{改正}~平台漏标(本卷补标)\quad
  {\color{cFix}\rule{1em}{0.6em}}~范文/正确\quad
  {\color{cNote}\rule{1em}{0.6em}}~AI 建议
  \end{tcolorbox}}

% ---------- 页眉页脚 ----------
\pagestyle{fancy}
\fancyhf{}
\renewcommand{\headrulewidth}{0.4pt}
\renewcommand{\footrulewidth}{0pt}
\fancyhead[L]{\footnotesize\color{cHead}\bfseries PTE Core 模考 1B}
\fancyhead[R]{\footnotesize\color{gray} 错题与详解整理卷}
\fancyfoot[C]{\footnotesize\color{gray}\thepage}

\begin{document}

% --------- 封面（原 cover.tex）---------
% ---------- 封面 / 成绩总览 ----------
\definecolor{mListen}{HTML}{1F3A5F}
\definecolor{mRead}{HTML}{C2D70B}
\definecolor{mSpeak}{HTML}{4D4D4D}
\definecolor{mWrite}{HTML}{A01B7D}

\begin{titlepage}
\thispagestyle{empty}
\centering
\vspace*{1.4cm}

{\fontsize{30}{36}\selectfont\bfseries\color{cHead} PTE Core 模考 1B}\\[4mm]
{\fontsize{17}{20}\selectfont\color{cAccent}\bfseries 错题与详解 · 整理卷}\\[3mm]
{\color{gray} APEUni / 猩际 PTE 模考　·　提交时间 2026-06-21 12:36}\\[1.2cm]

% ---- 总分 ----
\begin{tcolorbox}[width=0.62\textwidth, halign=center, enhanced,
  colback=cAccent, colframe=cAccent, arc=6pt, boxrule=0pt, top=10pt, bottom=10pt]
  {\color{white}\large 总分 Overall}\\[2pt]
  {\fontsize{46}{50}\selectfont\bfseries\color{white} 40}
\end{tcolorbox}

\vspace{1.0cm}

% ---- 四项分数条形图 ----
\begin{tcolorbox}[width=0.84\textwidth, enhanced, colback=white, colframe=cRule,
  boxrule=0.8pt, arc=4pt, left=18pt, right=18pt, top=14pt, bottom=14pt]
\setlength{\tabcolsep}{6pt}\renewcommand{\arraystretch}{1.7}
\sffamily
\begin{tabular}{@{}r@{\hskip 8pt}l@{\hskip 10pt}l@{}}
 \bfseries Listening 听力 & {\color{mListen}\rule{3.67cm}{0.42cm}} & {\bfseries\large 33}\\
 \bfseries Reading 阅读   & {\color{mRead}\rule{4.67cm}{0.42cm}}   & {\bfseries\large 42}\\
 \bfseries Speaking 口语  & {\color{mSpeak}\rule{3.22cm}{0.42cm}}  & {\bfseries\large 29}\\
 \bfseries Writing 写作   & {\color{mWrite}\rule{6.33cm}{0.42cm}}  & {\bfseries\large 57}\\
\end{tabular}\\[2mm]
{\footnotesize\color{gray} 满分 90　|　刻度: 条长 = 得分 / 90}
\end{tcolorbox}

\vfill
\begin{flushleft}\small\color{gray}
\textbf{说明:}本卷由本次模考的题目与"详解"截图逐题誊写、重新排版而成,完整保留每道题的
\emph{题干 / 原文、范文、你的作答、逐项评分、彩色批改与 AI 中文建议}。\\
所有错误按平台标注用颜色标出(见下方图例)。原始"详解"PDF 不可复制、仅截图,本卷内容均为人工对照誊录,若与原图有出入以原图为准。
\end{flushleft}
\vspace{4mm}
\ptelegend
\end{titlepage}

% --------- 听力正文（原 sec_listening_body.tex）---------
% ===================== 听力 Listening =====================
\ptesection{听力}{Listening · SST + MCM + FIB + MCS + SMW + HIW + WFD}{33}

\vspace{1mm}
{\small\color{gray} 依据「听力1 / 听力2」逐题誊录，并用「听力0 / 听力详解1」补全 SST 评分明细、用长截图补全 PDF 中被悬浮按钮遮挡处。
本套听力共 \textbf{15} 题：SST 概述口语$\times$1、MCM-L 多选$\times$1、FIB-L 听力填空$\times$3、MCS-L 单选$\times$1、SMW 选词$\times$2、HIW 找错词$\times$3、WFD 听写$\times$4。}
\vspace{1mm}

\begin{notebox}
{\footnotesize 标注沿用阅读：{\color{cFix}\ding{51}\,绿=对/命中}；{\color{cWrong}\ding{55}\,红删=错}，其后 {\color{cFix}(答案/听: 正确)}。
\textbf{HIW}（找错词）：屏幕上与录音不符的词用 {\color{cWrong}\uline{红下划线}} 标出（即应点选的"错词"），括注实际读音。
\textbf{WFD}（听写）：{\color{cFix}绿}=命中词，{\color{cWrong}\sout{红删}}=多写/听错的词。}
\end{notebox}

% =================== Q1 ===================
\begin{qblock}{Q1\quad SST (C) \#21\ \ Role of Education}{概述口语 Summarize Spoken Text · 评分 \textcolor{cAccent}{\bfseries 7 / 10}}
\ptesub{录音原文 Transcript}
\begin{promptbox}
Education stands as a powerful tool in bridging social inequalities, serving as a catalyst for change in a world riddled with disparities. Its role in shaping a more equitable society is paramount, as it offers the key to unlocking potential and leveling the playing field. However, the journey towards educational equity is fraught with challenges. In many parts of the world, access to quality education is unevenly distributed, often influenced by factors such as socioeconomic status, race, and gender. This disparity perpetuates a cycle of inequality, where underprivileged groups remain marginalized. The digital divide, exacerbated by the recent global pandemic, has further highlighted these disparities. Students in remote or impoverished areas, lacking access to technology and the internet, find themselves at a significant disadvantage. However, there is hope. Initiatives aimed at providing equitable access to education, such as scholarship programs, community schools, and digital learning platforms, have shown promise. They not only offer access to knowledge but also foster critical thinking and problem-solving skills, empowering individuals to rise above their circumstances.
\end{promptbox}
\zh{教育是弥合社会不平等的有力工具，在一个充满差距的世界里充当变革的催化剂。它在塑造更公平社会中的作用至关重要，因为它提供了释放潜能、创造公平竞争环境的钥匙。然而，通往教育公平的道路充满挑战。在世界许多地方，优质教育的获取分布不均，常受社会经济地位、种族和性别等因素影响。这种差距使不平等的循环延续，弱势群体始终被边缘化。近期的全球疫情加剧了数字鸿沟，进一步凸显了这些差距：身处偏远或贫困地区、缺乏技术与网络的学生处于明显劣势。然而，希望仍在。旨在提供公平教育机会的举措——如奖学金项目、社区学校和数字学习平台——已显现成效；它们不仅提供获取知识的途径，还培养批判性思维和解决问题的能力，使个人能够超越自身处境。}
\ptesub{你的概述 Your Summary}
\begin{promptbox}
The education is important. {\color{cAccent}\uline{The society}} has some {\color{cAccent}\uline{journal}} challenges. The influence {\color{cAccent}\uline{contain}} society and gender. The community has some critical skills.
\end{promptbox}
{\footnotesize\color{gray}（AI 标注待改：The/The society 多余冠词；journal 疑为听错词（general?）；contain → contains 主谓一致。）}
\ptesub{AI 评分明细（本题 10 分，得 \textcolor{cAccent}{\bfseries 7} 分）}
\begin{notebox}\small
\begin{tabular}{@{}lcl@{}}
Content（内容）& \textbf{2 / 2} & 内容很棒\\
Form（字数形式）& \textbf{2 / 2} & 很棒，这是必拿分数\\
Grammar（语法）& \textcolor{cWrong}{\bfseries 0 / 2} & SST 中占分比重、却很容易提高；丢分较多，建议针对性练习\\
Spelling（拼写）& \textbf{2 / 2} & 拼写很棒，再接再厉\\
Vocabulary（词汇）& \textbf{1 / 2} & PTE 不要求难词，但需用词恰当\\
\end{tabular}
\end{notebox}
\end{qblock}

% =================== Q2 ===================
\begin{qblock}{Q2\quad MCM-L \#20\ \ Re-adjusting to Life}{多选 Multiple Answers · 评分 \textcolor{cAccent}{\bfseries 1 / 2}}
\ptesub{录音原文 Transcript}
\begin{promptbox}
Re-adjusting to life in your own country after living abroad for some time is a little like recovering from jet-lag after a long flight across several time zones. It takes time. And research indicates that after nine years living in a foreign country you never really do re-adjust. Of course, things have changed: governments have come and gone, what you knew as countryside has become a suburb, new technologies have changed the way people go about their daily lives, and so on. These changes may well have been taking place in your adopted country, but they were happening while you were there, so you could adapt as you went along. Those are not the main difficulties, however. It is in the smaller, everyday things that you might experience what is known as culture shock, although it's not really a shock, but puzzling all the same. For example, the precise way to behave at a supermarket checkout may have changed. And in ordinary conversation, the frames of reference have changed, and quite often you find that you don't really know what people are talking about, even though they are speaking your native tongue.
\end{promptbox}
\zh{在国外生活一段时间后回到自己国家重新适应，有点像长途跨时区飞行后倒时差——需要时间。研究表明，在国外生活九年后，你其实永远无法真正重新适应。当然，一切都变了：政府换了一届又一届，你记忆中的乡村变成了郊区，新技术改变了人们的日常生活方式，等等。这些变化在你侨居的国家或许也在发生，但那是在你在场时发生的，所以你能边走边适应。然而，那些并不是主要困难。真正的困难在于更小的日常琐事——你可能经历所谓“文化冲击”，虽算不上冲击，却同样令人困惑。比如超市结账时的确切做法可能变了；日常对话中参照框架也变了，你常常发现自己根本不知道别人在说什么，尽管他们说的是你的母语。}
\ptesub{题目 Question}
\begin{notebox}\small Which of the following are mentioned as being difficult to re-adapt to when returning to your native country after a long absence?\end{notebox}
\ptesub{选项 Choices（多选）}
{\small
A）Technological change\\
B）More than nine years absence\\
C）The difference in time zones\\
{\color{cFix}\textbf{D）The way people go about their daily lives}}\\
{\color{cFix}\textbf{E）People's terms of reference in conversation}}\\
F）Changes in government}
\smallskip\par
{\small 正确答案 \textcolor{cFix}{\bfseries D, E}\qquad 你的答案 {\color{cWrong}\bfseries A,} {\color{cFix}\bfseries D, E}\;{\footnotesize\color{gray}（D、E 对，多选 A 错扣分）}}
\end{qblock}

% =================== Q3 ===================
\begin{qblock}{Q3\quad FIB-L \#155\ \ Transport Chaos}{听力填空 Fill in the Blanks · 评分 \textcolor{cAccent}{\bfseries 1 / 4}}
\ptesub{录音原文 Transcript（含填空作答）}
\begin{promptbox}
The New South Wales government has \rok{apologized} for yesterday's transport chaos in and around Sydney Harbour during the visit of the Queen Mary II and the Queen Elizabeth II. Roads were jammed, traffic ground to a halt while tram and \rno{serious}{ferry} services were swamped with thousands of additional passengers, with most services delayed for hours. Premier Maurice Humor says that plans were put in place to ``deal with the congestion'' but the number of visitors well exceeded \rno{expantation}{expectation}. On the harbour itself there seemed to be as much congestion as there was on the roads, but everyone agreed it was an amazing \rno{spectical}{spectacle}.
\end{promptbox}
\zh{新南威尔士州政府就昨天玛丽女王二号与伊丽莎白女王二号到访期间、悉尼港周边的交通混乱致歉。道路被堵死、交通陷入停滞，有轨电车与渡轮服务被成千上万的额外乘客挤爆，大多数班次延误数小时。州长 Maurice Humor 称已制定方案“应对拥堵”，但访客人数远超预期。港口本身的拥堵程度似乎与道路不相上下，但人人都认同这是一场惊人的盛况。}
\par{\footnotesize\color{gray}（FIB-L 为听写打字，无选项；你的作答：apologized / serious / expantation / spectical —— 后三处分别听错词或拼写错误。）}
\end{qblock}

% =================== Q4 ===================
\begin{qblock}{Q4\quad FIB-L \#164\ \ Dogs}{听力填空 Fill in the Blanks · 评分 \textcolor{cAccent}{\bfseries 0 / 3}}
\ptesub{录音原文 Transcript（含填空作答）}
\begin{promptbox}
Dogs are not just man's best friend. Previous studies have shown that kids with dogs are less likely to develop asthma. Now a new study may show how --- if results from mice apply to us. The work was presented at a meeting of the American Society for Microbiology. The study tests what's called the \rno{highjeans}{hygiene} hypothesis. The idea is that extreme cleanliness may actually promote disease later on. Researchers collected dust from homes that had a dog. They fed that house dust to mice. They then infected the mice with a common \rno{im}{childhood} infection called \rno{rest}{respiratory} syncytial virus --- or RSV.
\end{promptbox}
\zh{狗不只是人类最好的朋友。以往研究表明，养狗家庭的孩子更不易患哮喘。如今一项新研究或许能揭示其原理——前提是小鼠实验结果同样适用于我们。研究在美国微生物学会的一次会议上发表，检验的是所谓“卫生假说”：过度清洁反而可能在日后诱发疾病。研究者从养狗家庭收集灰尘喂给小鼠，再让小鼠感染一种常见的儿童感染病——呼吸道合胞病毒（RSV）。}
\par{\footnotesize\color{gray}（你的作答：highjeans / im / rest —— 三处均听错。）}
\end{qblock}

% =================== Q5 ===================
\begin{qblock}{Q5\quad FIB-L \#134\ \ Almonds}{听力填空 Fill in the Blanks · 评分 \textcolor{cAccent}{\bfseries 1 / 4}}
\ptesub{录音原文 Transcript（含填空作答）}
\begin{promptbox}
And one particular crop, almond in the US and now in Australia, is \rok{transforming} the world of beekeeping and of bees. What has happened is that something serendipitous came along that people found out, that doctors found out that almonds are good for you, a \rno{confaction}{confection} but it's good for you. The Almond Board got a very aggressive promotion going on for almonds. They actually, I just heard recently, sent out sales reps to \rno{cardiology}{cardiologists} at hospitals to promote the heart benefits of almonds. In a very good promotion of almonds, and it's a \rno{magenal}{legitimate} promotion because they are a healthy food.
\end{promptbox}
\zh{在美国、如今也在澳大利亚，有一种特别的作物——杏仁——正在改变养蜂业与蜜蜂的世界。事情是这样的：人们、医生偶然发现杏仁对健康有益——它虽是一种甜食，却对你有好处。杏仁委员会为杏仁开展了非常激进的推广。我最近才听说，他们甚至派销售代表去医院找心脏科医生，宣传杏仁对心脏的益处。这是对杏仁很好的推广，而且是正当的推广，因为它们确实是健康食品。}
\par{\footnotesize\color{gray}（你的作答：transforming / confaction / cardiology / magenal —— 仅首词正确，其余听错或拼写错误。）}
\end{qblock}

% =================== Q6 ===================
\begin{qblock}{Q6\quad MCS-L \#47\ \ Handball Violation}{单选 Single Answer · 评分 \textcolor{cAccent}{\bfseries 1 / 1}}
\ptesub{录音原文 Transcript}
\begin{promptbox}
Listen to part of a lecture and answer the question. Ok, now let's talk about handball violations. I bet you all know that you aren't supposed to use your hands when you play soccer. It is called a handball violation. You cannot touch any part of the ball with your hands unless you are the goalkeeper. Only the goalkeeper can touch the ball with his hands. A handball violation includes using any part of the body from the fingers to the shoulders. Until here. This is all you know, now I want to ask you some questions. First, picture the situation in the middle of a game. A player kicks the ball and the ball touches the hand of the opposite team player, but he didn't touch it intentionally. Now, is it a handball violation? I want you to think about it. Here goes another question. What if the kicker did that intentionally? Is it a handball violation?
\end{promptbox}
\zh{听一段讲座并回答问题。好，现在我们来谈谈手球犯规。我打赌你们都知道，踢足球时不能用手——这叫手球犯规。除非你是守门员，否则不能用手触球的任何部位；只有守门员能用手触球。手球犯规包括用从手指到肩膀的身体任何部位触球。先讲到这。以上都是你们已知的，现在我要问几个问题。首先，设想比赛进行中的情形：一名球员踢球，球碰到了对方球员的手，但他不是故意触球的——那么，这算手球犯规吗？我想让你们想一想。再来一个问题：如果是踢球者故意这么做的呢？那算手球犯规吗？}
\ptesub{题目 Question}
\begin{notebox}\small What will follow this lecture?\end{notebox}
\ptesub{选项 Choices（单选）}
{\small
{\color{cFix}\textbf{A）The role of intentions in handball violations}}\\
B）Handball violations by goalkeepers\\
C）The referee in soccer games\\
D）A soccer game without rules}
\smallskip\par
{\small 正确答案 \textcolor{cFix}{\bfseries A}\qquad 你的答案 \textcolor{cFix}{\bfseries A}\;{\footnotesize\color{gray}（讲座最后连续抛出“是否故意”的问题，下文将讨论“意图”的作用）}}
\end{qblock}

% =================== Q7 ===================
\begin{qblock}{Q7\quad SMW \#51\ \ Migration}{选词 Select Missing Word · 评分 \textcolor{cAccent}{\bfseries 1 / 1}}
\ptesub{录音原文 Transcript（结尾词被“哔”声替换）}
\begin{promptbox}
In the nineteenth century, there were several periods when large numbers of people moved from one place to another around the world. In many cases, people move to another continent. These mass migrations were on a much larger scale than any previous migrations in history. One major movement was from Europe to the Americas, Australia and Africa. This migration of Europeans involved around sixty million people over one hundred years. Another mass migration was from Russia to Siberia and central Asia. Another was from China, India and Japan to Southeast Asia. These large movements of people were made possible by the new cheap and fast means of \rule{1.2cm}{0.4pt}\ \textit{[beep]}.
\end{promptbox}
\zh{19 世纪有几个时期，世界各地大量人口从一地迁往另一地，很多情况下迁往另一个大洲。这些大规模迁徙的规模远超历史上任何以往的迁徙。一次主要迁移是从欧洲到美洲、澳大利亚和非洲，这场欧洲人的迁徙在一百年间涉及约六千万人；另一次是从俄国到西伯利亚和中亚；还有一次是从中国、印度、日本到东南亚。这些大规模的人口流动，是靠新型廉价而快速的【交通工具】才得以实现的。}
\ptesub{选项 Choices（单选）}
{\small A）transition\quad B）internet\quad {\color{cFix}\textbf{C）transportation}}\quad D）airlines}
\smallskip\par
{\small 正确答案 \textcolor{cFix}{\bfseries C}\qquad 你的答案 \textcolor{cFix}{\bfseries C}\;{\footnotesize\color{gray}（cheap and fast means of \emph{transportation}）}}
\end{qblock}

% =================== Q8 ===================
\begin{qblock}{Q8\quad SMW \#48\ \ Worker Bee}{选词 Select Missing Word · 评分 \textcolor{cAccent}{\bfseries 0 / 1}}
\ptesub{录音原文 Transcript（结尾词被“哔”声替换）}
\begin{promptbox}
The body of the worker bee is divided into three segments, head, thorax, and abdomen. On the head are the mandible, antennae, jar-like organs, which enable the bees to perform the necessary high duties and to mold the wax and build their combs. The honey bee's four wings and six legs are fastened to the thorax. Located in the abdomen are the honey sac and the sting with its highly developed poison sac. The sting is used by the workers for self defense and for the protection of their colony. The worker uses her sting only once, for in doing so she \rule{1.2cm}{0.4pt}\ \textit{[beep]}.
\end{promptbox}
\zh{工蜂的身体分为三节：头、胸、腹。头部有上颚、触角等罐状器官，使蜜蜂能完成必要的高强度工作，并塑造蜂蜡、筑造蜂巢。蜜蜂的四翅六足连在胸部。腹部有蜜囊以及带高度发达毒囊的螫针。螫针被工蜂用于自卫和保护蜂群。工蜂的螫针只能用一次，因为这样做时她就【失去了生命】……}
\ptesub{选项 Choices（单选）}
{\small A）rescues herself\quad B）loses her dignity\quad {\color{cFix}\textbf{C）loses her life}}\quad D）kills the enemy}
\smallskip\par
{\small 正确答案 \textcolor{cFix}{\bfseries C}\qquad 你的答案 {\color{cWrong}\bfseries A}\;{\footnotesize\color{gray}（蜜蜂蜇人后螫针留在体内，故“失去生命”）}}
\end{qblock}

% =================== Q9 ===================
\begin{qblock}{Q9\quad HIW \#70\ \ Diabetes}{找错词 Highlight Incorrect Words · 评分 \textcolor{cAccent}{\bfseries 1 / 6}}
\ptesub{录音原文 Transcript（屏幕错词已标出，括注实际读音）}
\begin{promptbox}
No that was, and that's an important aspect, as you \hiw{referred}{alluded} to earlier we've previously done work which has proven that in some \hiw{circumstances}{situations}, even people whose blood pressure is not high, can benefit from blood pressure lowering \hiw{rehabilitation}{therapy}. So in this study the main reason that we included the patients was because of diabetes, we didn't care what their blood pressure was, whether it was high or low. And our \hiw{intention}{objective} was to not lowering average or below average blood pressure in diabetics was beneficial and the \hiw{effect}{result} suggested that irrespective of whether your blood pressure was high or low, if you had diabetes you \hiw{profited}{benefited}.
\end{promptbox}
\zh{不，那是——这是个重要方面——正如你前面提到的，我们之前的研究已证明：在某些情况下，即便血压不高的人，也能从降血压治疗中获益。在这项研究里，我们纳入患者的主要原因是糖尿病，并不在意他们血压是高是低。我们的目标是考察——把糖尿病患者本就平均或偏低的血压再降低是否有益；结果表明，无论血压高低，只要患有糖尿病，都能获益。}
\par{\small 应点的 6 个错词：{\color{cWrong}referred / circumstances / rehabilitation / intention / effect / profited}\quad（实际读音 alluded / situations / therapy / objective / result / benefited）。\;你只点了 {\color{cWrong}referred}。}
\end{qblock}

% =================== Q10 ===================
\begin{qblock}{Q10\quad HIW \#317\ \ Classified Advertisements}{找错词 Highlight Incorrect Words · 评分 \textcolor{cAccent}{\bfseries 1 / 6}}
\ptesub{录音原文 Transcript（屏幕错词已标出，括注实际读音）}
\begin{promptbox}
Classified advertisements placed by individuals in \hiw{newsprint}{newspapers} and magazines are not covered by the Advertising Standards Authority's `\hiw{court}{code} of practice'. If you happen to buy goods that have been wrongly described in such an advertisement, and have lost money as a result, the only thing you can do is bring a case against the person who placed the advertisement for misrepresentation or for breach of \hiw{contrast}{contract}. In this case, you would use the small claims procedure, which is a relatively cheap way to sue for the recovery of a debt. If you want to pursue a claim, you should take into account whether the person you are suing will be able to pay damages, should any be \hiw{rewarded}{awarded}. Dishonest traders are \hiw{wary}{aware} of this and often pose as private sellers to \hiw{expose}{exploit} the legal loopholes that exist: that is, they may claim they are not in a position to pay damages.
\end{promptbox}
\zh{个人在报纸和杂志上刊登的分类广告，不受广告标准管理局《行为准则》的约束。如果你碰巧买到在此类广告中被错误描述的商品并因此蒙受损失，你唯一能做的就是以虚假陈述或违约为由起诉刊登广告的人。这种情况下你会走小额索赔程序——一种相对便宜的追讨欠款诉讼方式。若要索赔，应考虑被告是否有能力支付应判的赔偿金。不诚实的商家对此心知肚明，常假扮成私人卖家来利用现有法律漏洞：即声称自己无力支付赔偿。}
\par{\small 应点的 6 个错词：{\color{cWrong}newsprint / court / contrast / rewarded / wary / expose}\quad（实际读音 newspapers / code / contract / awarded / aware / exploit）。\;你只点了 {\color{cWrong}newsprint}。}
\end{qblock}

% =================== Q11 ===================
\begin{qblock}{Q11\quad HIW \#316\ \ to Travel Abroad}{找错词 Highlight Incorrect Words · 评分 \textcolor{cAccent}{\bfseries 3 / 6}}
\ptesub{录音原文 Transcript（屏幕错词已标出，括注实际读音）}
\begin{promptbox}
In the 19th century, few people could afford to travel abroad; it was expensive and there weren't the \hiw{massive}{mass} transport systems that we have today. So curiosity about foreign lands had to be satisfied through books and drawings. With the advent of photography, a whole new \hiw{version}{dimension} of `reality' became available. Publishers were not slow to realize that here was a large new market of people \hiw{eager}{hungry} for travel photography and they soon had photographers out shooting the best known European cities, as well as more exotic places further \hiw{afield}{away}. People bought the pictures by the millions, and magic lantern shows were presented in schools and \hiw{leisure}{lecture} halls. Most popular of all, however, was the stereoscopic picture which \hiw{pretended}{presented} three-dimensional views and was considered a marvel of Victorian technology.
\end{promptbox}
\zh{19 世纪很少有人负担得起出国旅行——费用昂贵，也没有我们今天这样庞大的交通系统。因此人们对异国的好奇只能靠书籍和绘画来满足。随着摄影术出现，一种全新维度的“现实”变得可得。出版商很快意识到，这里有一个庞大的新市场——人们渴望旅行摄影——于是很快派摄影师去拍摄最著名的欧洲城市以及更遥远的异域。人们成百万地购买这些照片，学校和讲堂也放映幻灯片表演。然而最受欢迎的是立体照片，它呈现三维视图，被视为维多利亚时代技术的奇迹。}
\par{\small 应点的 6 个错词：{\color{cWrong}massive / version / eager / afield / leisure / pretended}\quad（实际读音 mass / dimension / hungry / away / lecture / presented）。\;你点了 {\color{cFix}eager, afield, pretended}（3 个全对）。}
\end{qblock}

% =================== Q12 ===================
\begin{qblock}{Q12\quad WFD \#1667}{听写句子 Write From Dictation · 评分 \textcolor{cAccent}{\bfseries 3 / 8}}
\ptesub{正确句子 Correct Sentence}
\begin{promptbox}
Different factors affect the freezing time of water.
\end{promptbox}
\zh{不同的因素会影响水的结冰时间。}
\ptesub{你的作答（平台逐词批改：{\color{cFix}绿=命中}，{\color{gray}(灰)=漏写}，{\color{cWrong}红删=多写/听错}）}
{\small\lhit{Different} {\color{gray}(factors affect the freezing)} \lhit{water} \lmiss{often} \lhit{time} {\color{gray}(of)} \lmiss{times}}
\end{qblock}

% =================== Q13 ===================
\begin{qblock}{Q13\quad WFD \#1666}{听写句子 Write From Dictation · 评分 \textcolor{cAccent}{\bfseries 6 / 8}}
\ptesub{正确句子 Correct Sentence}
\begin{promptbox}
The subject is complex and difficult to explain.
\end{promptbox}
\zh{这个主题复杂且难以解释。}
\ptesub{你的作答（{\color{cFix}绿=命中}，{\color{gray}(灰)=漏写}，{\color{cWrong}红删=多写/听错}）}
{\small\lhit{The subject is} {\color{gray}(complex and)} \lmiss{somtimes} \lhit{difficult to explain}}
\end{qblock}

% =================== Q14 ===================
\begin{qblock}{Q14\quad WFD \#1665}{听写句子 Write From Dictation · 评分 \textcolor{cAccent}{\bfseries 1 / 9}}
\ptesub{正确句子 Correct Sentence}
\begin{promptbox}
Experience would be an advantage for this managerial role.
\end{promptbox}
\zh{经验对这个管理职位将是一项优势。}
\ptesub{你的作答（{\color{cFix}绿=命中}，{\color{gray}(灰)=漏写}，{\color{cWrong}红删=多写/听错}）}
{\small{\color{gray}(Experience would be an advantage)} \lmiss{Explosions is} \lhit{for} {\color{gray}(this managerial role)} \lmiss{loyal raw lay}}
\end{qblock}

% =================== Q15 ===================
\begin{qblock}{Q15\quad WFD \#1663}{听写句子 Write From Dictation · 评分 \textcolor{cAccent}{\bfseries 5 / 9}}
\ptesub{正确句子 Correct Sentence}
\begin{promptbox}
Background music can help students concentrate on their studies.
\end{promptbox}
\zh{背景音乐能帮助学生集中精力学习。}
\ptesub{你的作答（{\color{cFix}绿=命中}，{\color{gray}(灰)=漏写}，{\color{cWrong}红删=多写/听错}）}
{\small{\color{gray}(Background)} \lmiss{Backgound} \lhit{music can help} {\color{gray}(students)} \lmiss{student make more} \lhit{concentrate} \lmiss{concentrations concentration} \lhit{on} {\color{gray}(their studies)} \lmiss{study}}
\end{qblock}

\end{document}
